\section{The rotated threefold}\label{sec:rotated}

Choose the circle with character \(t\) as the external circle. Its fixed
threefold is \(X'=C\times S\), where \(S=\mathbb A^2\) carries the two
weights \(q,q^{-1}\). A curve can now move off the origin of \(S\) by a
constant translation. We compute on this noncompact moduli space before
pushing its index to cycles.

\subsection{Deforming the framed matrix equation}

Let \(\mathcal O_{X'}\to F\) be a stable pair of curve degree \(m>0\) and
Euler characteristic \(n\). Its proper support projects finitely to \(C\).
Indeed, the projection is affine and its restriction to the support is
proper. Purity then makes \(E=p_*F\) a vector bundle of rank \(m\),
where \(p:X'\to C\). Riemann--Roch gives
\begin{equation}
 k=\deg E=n+m(g-1).
 \label{eq:rotated-excess}
\end{equation}
The two coordinates of \(S\) act on \(E\), and the section of the pair
gives a distinguished vector \(i:\mathcal O_C\to E\). Over the function
field of \(C\), the vector generates \(E\) under the coordinate actions.

The framed ADHM description of the stable-pair theory adds a covector
\(j:E\to\mathcal O_C\) and writes the relation as a matrix equation
\cite[Theorem~1.9 and \S7]{diaconescu}. Write \(x,y\) for the two
endomorphisms. We deform the equation, with its scalar part included, to
\begin{equation}
 [x,y]+ij=a\,1_E,\qquad a\in\mathbb A^1 .
 \label{eq:rotated-adhm}
\end{equation}
Framing fixes the copy of \(\mathcal O_C\), and we divide by changes of
basis of \(E\). The stable open is specified by generic generation by
the words in \(x,y\) applied to \(i\). At \(a=0\) this is
Diaconescu's presentation of the stable-pair moduli. Under the opposite
weights on \(x,y\), the parameter \(a\) is invariant.

At \(a\ne0\), fix any point of \(C\). The span of the words applied to
\(i\) is invariant under \(x,y\). On its quotient the commutator would be
\(a\) times the identity. Taking the trace forces this quotient to vanish.
Dually, the common kernel of the covectors \(j\,w(x,y)\), for all words
\(w\), is invariant and contained in \(\ker j\). The same trace argument
on this subspace makes it zero. Words of length at most \(m-1\) suffice
in each argument, since an increasing span which stops growing is already
invariant. These words are global sections on \(C\). Consequently \(E\)
and \(E^\vee\) are both globally generated and \(\deg E=0\). At \(a=0\)
there are still \(m\) words generically independent on \(i\). Their
determinant is a nonzero section of \(\det E\), so \(k\ge0\).

For \(k>0\) this calculation implies the integral identity
\begin{equation}
 [\mathcal O^{\mathrm{vir}}_{P_n(X',m)}]=0
 \qquad\text{in }G_0^{\mathbb C^\times_q}(P_n(X',m)).
 \label{eq:rotated-vanishing}
\end{equation}
We give the specialization argument, since the equivariant weight of
\(a\) matters here. Use the derived zero locus of
\eqref{eq:rotated-adhm} modulo changes of frame and take its stable space
of sections over the fixed curve. Relative to the \(a\)-line its tangent
complex is the derived push from \(C\) of the complex in degrees
\(-1,0,1\)
\begin{equation}
 \operatorname{End}E
 \longrightarrow (\operatorname{End}E)q\oplus
  (\operatorname{End}E)q^{-1}\oplus E\oplus E^\vee
 \longrightarrow \operatorname{End}E .
 \label{eq:rotated-deformation}
\end{equation}
The first arrow is the infinitesimal change of basis and the second is
the differential of the moment map. Generic cyclicity makes the first
injective over the function field. The trace symplectic pairing on the
framed matrix space identifies its transpose, up to sign, with the
second arrow, which is therefore generically surjective. On \(C\) the
first kernel sheaf vanishes and the last cokernel has finite support.
Pushing to the parameter space gives tangent amplitude \([0,1]\).
This is the fixed-domain quasimap obstruction theory
\cite[\S4.3]{okounkov}. Stabilizers are trivial by generic cyclicity.
Boundedness at fixed \(k\) follows by dualizing \(m\) generically
independent word sections. The resulting subbundle away from finitely
many points, \(E^\vee\subset\mathcal O_C^m\), has quotient of length
\(k\), hence belongs to a finite-type Quot scheme.

For clarity, the comparison on the central fiber can be made at the
level of virtual structure sheaves. If \(\pi\) forgets the \(C\)-factor
of the universal family, the stable-pair virtual tangent has class
\begin{equation}
 R\pi_*\!\left(E+E^\vee-
       (\operatorname{End}E)(1-q)(1-q^{-1})\right).
 \label{eq:rotated-tangent}
\end{equation}
The Koszul resolution for the two plane coordinates gives the quadratic
term in \(E\) in the trace-free pair complex. The two linear terms come
from the section and its dual, with the product of the two plane weights
equal to one. This agrees with the alternating class of
\eqref{eq:rotated-deformation} on the same moduli scheme. The
\(K\)-theoretic Fulton class formula expresses the virtual structure
sheaf of a quasi-projective scheme in terms of its intrinsic Fulton
class and its virtual cotangent class
\cite[Theorem~4.2 and Corollary~4.5]{thomas}. Its cone filtration and
exterior powers preserve the \(q\)-linearizations, with a separate
auxiliary grading for the filtration. Thus the comparison holds in
integral equivariant coherent \(K\)-theory. See also the local-curve
application in \cite[\S9]{monavari}.

The virtual structure sheaf on the central fiber is the refined pullback
of the relative virtual family along \(a=0\), by derived base change
or virtual specialization \cite{qu}. For \(k>0\) the underlying space
of the entire family lies over zero. Filter a coherent sheaf there by
powers of \(a\). Every successive quotient is a sheaf on the central
fiber, and its refined pullback has two equal Tor terms in adjacent
degrees, since \(a\) has weight zero. The pullback thus kills every
class of the total family. This proves \eqref{eq:rotated-vanishing},
even when the parameter acts nilpotently on its ordinary structure
sheaf.

\subsection{The plane Hilbert scheme and its orientation}

At \(k=0\) the determinant of the \(m\) independent word sections has
degree zero and no zero. These sections trivialize \(E\) on \(C\).
The coordinate operators and the cyclic vector of the stable pair are
then constant, and give a cyclic length-\(m\) quotient of the coordinate
ring of \(S\). This identification works in families. Locally on the
parameter scheme choose words that form a basis on one fiber over a
point of \(C\). At the chosen parameter they form a basis on the entire
curve. Properness of \(C\) allows one to shrink the parameter scheme
until they form a basis everywhere. Since \(C\) is projective and
connected, the matrix entries then descend to the parameter scheme,
also over nonreduced bases. It follows that
\begin{equation}
 P_{-mb}(X',m)\simeq H_m=\operatorname{Hilb}^m(S),
 \qquad
 \xi\longmapsto [\mathcal O_{X'}\longrightarrow
                          \mathcal O_{C\times\xi}] .
 \label{eq:rotated-hilbert}
\end{equation}

For these constant cyclic data the framed complex on \(C\) is pulled
back from the tangent bundle \(T H_m\). The virtual tangent is
\(R\Gamma(C,\mathcal O_C)\otimes T H_m\). Its degree-zero cohomology is
the ordinary tangent bundle of the smooth Hilbert scheme. Its
obstruction bundle is \(H^1(C,\mathcal O_C)\otimes T H_m\). The
symplectic form on \(H_m\) is invariant under the antidiagonal circle
and trivializes the tangent determinant. We use the resulting relative
orientation, suppressing fixed one-dimensional vector spaces from \(C\),
which have trivial circle character.

There is also a determinant contribution from the two external lines
\cite{no}. Let \(V\) be the rank-\(m\) tautological bundle on \(H_m\),
let \(\pi_B:C\times B\to B\), and put \(A=R\pi_{B*}\mathcal L\).
Relative duality gives \(R\pi_{B*}\mathcal M=A^\vee[-1]\).
The unweighted ratio of the two determinants on \(H_m\times B\) is
\begin{equation}
 D=\frac{\det(V\otimes A)}{\det(V\otimes A^\vee[-1])}.
 \label{eq:rotated-external-determinant}
\end{equation}
Here both complexes on \(B\) have virtual rank zero and the same
determinant. Thus \(D\) is trivial and we take its trivial square root.
This argument uses perfect complexes throughout, including where the
individual spaces of sections of \(\mathcal L\) jump. The external
prefactor is \((-t)^{n+mb}\), hence is one on \(H_m\). The modified
class on the minimum component is therefore the pullback from \(H_m\)
of
\begin{equation}
 \mathcal F_m=\lambda_{-1}(T^\vee H_m)^g,\qquad
 \lambda_{-1}(U)=\sum_i(-1)^i[\wedge^i U].
 \label{eq:rotated-exterior}
\end{equation}

The effect of balance can be tested on an actual family of coincident
points. With \(B=\operatorname{Pic}^d C\) set \(\delta=d-(g-1)\).
The two ranks in \eqref{eq:rotated-external-determinant} become
\(\delta,-\delta\), whereas their determinants still agree. Thus
\(D=(\det V)^{2\delta}\) with oriented root \((\det V)^\delta\).
The punctual Hilbert square is a \(\mathbb P^1\) of one-dimensional
quotients of the cotangent plane at the origin. On this line the
length-two algebra has the constants and the universal quotient
\(\mathcal O(1)\) as its two graded pieces. The degree of the root on
this collision family is \(\delta\). The choice \(d=g-1\) is therefore
detected by the disappearance of the twist on the Hilbert scheme
itself.

\subsection{Comparison on the space of cycles}

The \(q\)-fixed ideals of \(H_m\) are the monomial ideals
\(\xi_\lambda\) for partitions \(\lambda\) of \(m\). To see that the
one-parameter subgroup has no further fixed components, first observe
that every fixed zero-cycle of \(S\) is concentrated at the origin.
The punctual Hilbert scheme is projective, so every connected component
of the fixed locus contains a point fixed by the full coordinate torus.
Those points are monomial ideals. At \(\xi_\lambda\) the tangent weights
restrict to
\begin{equation}
 \big\{q^{h(s)},q^{-h(s)}:s\in\lambda\big\},
 \label{eq:rotated-hook-weights}
\end{equation}
by the arm and leg formula \cite[Lemma~5.4.5]{haiman}. They are all nontrivial, so
each such point is isolated and reduced in the fixed locus.

The infinitesimal \(q\)-rotation is an equivariant section of \(T H_m\).
Its linearization at a fixed ideal is invertible by
\eqref{eq:rotated-hook-weights}. Its zeros are therefore simple and its
Koszul complex resolves their structure sheaf. Since \(g\ge1\), one
exterior Euler factor in \eqref{eq:rotated-exterior} resolves the fixed
points. The other \(g-1\) factors restrict to them. This proves the
identity
\begin{equation}
 \mathcal F_m
   =\sum_{|\lambda|=m}[\mathcal O_{\xi_\lambda}]
        \prod_{s\in\lambda}
            \big((1-q^{h(s)})(1-q^{-h(s)})\big)^{g-1}
 \quad\text{in }G_0^{\mathbb C^\times_q}(H_m).
 \label{eq:rotated-integral-concentration}
\end{equation}
In particular no character denominator is needed for the transverse
plane, even for an elliptic curve, when each point contributes one.

Let \(s_m:B\to W_m\) be the section of the common fivefold cycle space
which assigns the zero-section cycle \(m[C]\). Hilbert--Chow sends each
\(\xi_\lambda\) to the origin of \(\operatorname{Sym}^m S\) with
multiplicity \(m\), so its associated curve has cycle \(m[C]\).
These pushforwards are proper on the displayed representatives.
Equivalently one can embed both threefold cycle spaces in \(W_m\).
Compactification by compatible projective bundles makes containment
in each closed subbundle a closed condition on cycles, and one then
restricts to the open locus disjoint from infinity.

Writing \(Z_m(q)\) for \eqref{eq:hook}, the rotated pushforward is
\((s_m)_*(Z_m(q)\mathcal O_B)\) in \(G_0^{q,t}(W_m)\). In the original presentation
section~\ref{sec:original} gives the same class after summing the
Euler-characteristic grading with its external weight. This proves
the circle exchange of Theorem~\ref{thm:exchange}. One grading has
become an ordinary rotation of the plane, and its hook weights now
come from the tangent spaces to punctual ideals.
